\documentclass[aspectratio=169, 10pt]{beamer}
\usepackage{xeCJK} % 处理中文
\usepackage{fontspec}
\usepackage{cite}  % 引用管理
\usepackage{tikz}
\usetikzlibrary{positioning, shapes.geometric, arrows.meta, calc}
\usetheme{Madrid}
\usecolortheme{beaver} % 使用红色系主题，显得更严肃/醒目，也可改回 default

% 设置字体 (根据你的系统环境调整，Overleaf 上通常不需要额外配置，本地可能需要)
% \setCJKmainfont{SimSun}

% --- 标题信息 ---
\title{基于属性的 SD-WAN 控制系统}
% \subtitle{从底层协议封装到声明式属性驱动的演进}
\author[李梓勤]{李梓勤 \\ 指导教师：高凯}
\institute[SCU]{四川大学网络空间安全学院}
\date{2025 年 12 月 23 日}

% --- 自定义命令 ---
\newcommand{\citepeer}[1]{\textsuperscript{\cite{#1}}}

\begin{document}

% =================================================
% 1. 标题页
% =================================================
\frame{\titlepage}

% =================================================
% 2. 研究背景
% =================================================
\begin{frame}{背景}
    \begin{columns}
        \column{0.6\textwidth}
        \begin{block}{SD-WAN 的演进}
            \begin{itemize}
                \item \textbf{传统场景：} 聚焦于企业多分支机构与数据中心的高
                    效互联，解决专线成本高、调度难的问题
                    \cite{Jain2013, Hong2013, ONF2014}。
                \item \textbf{新兴趋势（个人/SOHO）：} 随着边缘计算与 HomeLab
                    的兴起，个人用户面临多设备（IoT、NAS、云服务器）跨地域、跨
                    运营商互联的刚需。
            \end{itemize}
        \end{block}

        \column{0.38\textwidth}
        \begin{alertblock}{零信任 (Zero Trust)}
            传统的边界防护已失效。部分 SD-WAN 引入了 NIST 的 ZTNA \cite{NIST2020} 理念：
            \begin{itemize}
                \item 从“基于位置的信任”转向 \textbf{“基于身份与属性的信任”}。
                \item 强调持续验证与最小权限。
            \end{itemize}
        \end{alertblock}
    \end{columns}

    \vspace{0.5em}
    % \textbf{核心矛盾：} 现有的企业级方案太重，消费级方案太黑盒，缺乏一种\textbf{轻量级、细粒度}的控制手段。
\end{frame}

% =================================================
% 3. 相关工作 (Data Plane & Control Plane)
% =================================================
\begin{frame}{相关工作}
    \textbf{1. 数据平面 (Data Plane)}
    \begin{itemize}
        \item \textbf{技术栈：} OpenVPN~\cite{OpenVPN}, IPsec~\cite{IPsec}, WireGuard~\cite{Donenfeld2017}。
        \item \textbf{痛点：} 提供了虚拟 L3 信道与加密原语，但“裸”协议配置极其复杂（Key Exchange, Route Table, Firewall Rules），缺乏集中管理能力。
    \end{itemize}

    \vspace{1em}
    \textbf{2. 控制平面 (Control Plane) 封装}
    \begin{table}[]
        \small
        \begin{tabular}{|l|l|l|}
            \hline
            \textbf{分类} & \textbf{代表工作} & \textbf{特点与局限} \\ \hline
            \textbf{端侧软件} & \begin{tabular}[c]{@{}l@{}}Tailscale
                \cite{Tailscale}, ZeroTier\cite{ZeroTier}, \\
                EasyTier\cite{EasyTier}, tinc \cite{tinc}\end{tabular} & \begin{tabular}[c]{@{}l@{}}基于 WireGuard 或自研协议；\\ \textbf{局限：} 策略表达能力固化，配置分散。\end{tabular} \\ \hline
            \textbf{网关硬件} & \begin{tabular}[c]{@{}l@{}}Oray (蒲公英) \cite{OrayPgy}, \\ GL.iNet (GoodCloud)\end{tabular} & \begin{tabular}[c]{@{}l@{}}集成度高，开箱即用；\\ \textbf{局限：} Vendor-Lock，难以扩展自定义逻辑。\end{tabular} \\ \hline
        \end{tabular}
    \end{table}
\end{frame}

% =================================================
% 4. 问题定义 (Motivation)
% =================================================
\begin{frame}{相关工作}
    尽管现有工具基本解决了“连通性”问题，但在“可控性”上仍存在本质缺陷：

    \begin{exampleblock}{缺陷 1：表达力不足}
        现有的 ACL 多基于静态的 Peer ID 或 IP 地址（如 \texttt{Allow 10.0.0.1 to 10.0.0.2}）。
        \textbf{缺失：} 无法表达高级安全语义，如“仅允许 \textit{打过最新补丁} 且 \textit{位于可信网络} 的设备访问核心数据库”。
    \end{exampleblock}

    \begin{exampleblock}{缺陷 2：配置碎片化}
        部分方案（如 tinc, 纯 WireGuard）需要在每台机器上手动维护 Peer 列表。
    \end{exampleblock}

    \begin{exampleblock}{缺陷 3：扩展性受限 (Limited Extensibility)。}
        难以集成动态路由逻辑。无法实现环境感知策略。
    \end{exampleblock}
\end{frame}

% =================================================
% 5. 拟解决思路 (Proposed Method)
% =================================================
\begin{frame}{系统设计}
    \textbf{逻辑集中化控制、声明式意图表达、高度可扩展插件架构}

    \begin{columns}[T]
        \column{0.55\textwidth}
        \begin{block}{自动化与扩展}
            \begin{itemize}
                \item \textbf{Zero-Touch 部署：} 节点凭 Token 接入，Agent 自动完成隧道建立、密钥交换与防火墙配置。
                \item \textbf{插件化架构：} 通过插件机制扩展控制平面功能（如：NAT 穿透）。
            \end{itemize}
        \end{block}

        \begin{block}{声明式策略引擎}
            \begin{itemize}
                \item \textbf{ABAC：} 基于身份、设备状态动态授权。
                    \cite{NISTABAC}\\
                \textit{\footnotesize 例：IF device.ispatched == True THEN Allow Access}
                \item \textbf{属性路由：} 根据网络属
                    性/流属性（延迟、成本、运营商）调整流量路径。
            \end{itemize}
        \end{block}

    \end{columns}
\end{frame}

\begin{frame}{系统架构}
    \centering

    % --- 第一层：逻辑控制中心 ---
    \begin{columns}[t]
        \column{0.8\textwidth}
        \begin{block}{控制平面 (Control Plane / "The Brain")}
            \centering
            \textbf{Central Controller} \\
            \vspace{0.2em}
            \footnotesize 身份认证 (Auth) \quad | \quad 属性数据库 (Attribute Store) \quad | \quad 策略引擎 (ABAC/ABR)
        \end{block}
    \end{columns}

    \vspace{0.5em}
    $\downarrow$ \textit{\scriptsize 声明式配置下发 (gRPC / TLS)} $\downarrow$
    % \vspace{0.5em}

    % --- 第二层：分布式节点 ---
    \begin{columns}[t]
        \column{0.45\textwidth}
        \begin{exampleblock}{边缘节点 A (Agent)}
            \footnotesize
            \textbf{属性上报：} OS, 延迟, 位置 \\
            \textbf{执行器：} WireGuard 接口管理
        \end{exampleblock}

        \column{0.45\textwidth}
        \begin{exampleblock}{边缘节点 B (Agent)}
            \footnotesize
            \textbf{属性上报：} 运营商, 负载状态 \\
            \textbf{执行器：} 路由表动态更新
        \end{exampleblock}
    \end{columns}

    \vspace{1em}

    % --- 第三层：数据平面 ---
    \begin{alertblock}{数据平面 (Data Plane)}
        \centering
        \textbf{P2P Mesh Network (Encrypted WireGuard Tunnels)} \\
        \footnotesize 流量不经过控制器，实现低延迟、高隐私的直接互联
    \end{alertblock}

\end{frame}

% =================================================
% 6. Roadmap
% =================================================
\begin{frame}{研究计划与 Roadmap}
    \begin{description}
        \item[2025.12 - 2026.01] \textbf{设计阶段：} 确定属性描述语言（DSL）规范，完成 Controller 原型设计。
        \item[2026.02] \textbf{核心开发：} 实现基于属性的 Policy Engine 及 WireGuard 接口封装。
        \item[2026.03] \textbf{扩展功能：} 开发 Agent 端环境感知探针（延迟/丢包检测）。
        \item[2026.04] \textbf{测试验证：} 在跨地域多云环境下部署，对比传统 Mesh 方案性能。
        \item[2026.05] \textbf{论文撰写：} 整理实验数据，完成毕业论文初稿。
        \item[2026.06] \textbf{答辩与交付：} 代码开源与毕业答辩。
    \end{description}
\end{frame}

% =================================================
% 7. 参考文献 (手动管理版，方便单文件编译)
% =================================================
\begin{frame}[allowframebreaks]{参考文献}
    % \tiny % 缩小字号以容纳更多引用
    \bibliographystyle{ieeetr}
    \begin{thebibliography}{99}

        % SD-WAN General
        \bibitem{Jain2013} Jain, S., Kumar, A., et al. "B4: Experience with a globally-deployed software defined WAN." \textit{SIGCOMM}, 2013.
        \bibitem{Hong2013} Hong, CY., Kandula, S., et al. "Achieving High Utilization with Software-Driven WAN" \textit{SIGCOMM}, 2013.
        \bibitem{ONF2014} Open Networking Foundation. "SDN architecture." \textit{ONF TR-502}, 2014.

        % Data Plane
        \bibitem{Donenfeld2017} Donenfeld, Jason A. "WireGuard: Next Generation Kernel Network Tunnel." \textit{NDSS}, 2017.
        \bibitem{OpenVPN} OpenVPN Inc. "OpenVPN Community Documentation." 2024.
        \bibitem{IPsec} Frankel, S., et al. "IP Security (IPsec) and Internet Key Exchange (IKE) Document Roadmap." \textit{RFC 6071}, 2011.

        % Zero Trust
        \bibitem{NIST2020} Rose, S., et al. "Zero Trust Architecture." \textit{NIST Special Publication 800-207}, 2020.

        % Existing Tools
        \bibitem{Tailscale} Tailscale Inc. "How Tailscale works." \textit{Tailscale Blog}, 2020.
        \bibitem{tinc} Guus Sliepen. "tinc VPN daemon." \textit{tinc-vpn.org}, 1998.
        \bibitem{Oray} Shanghai Oray Co., Ltd. "PgyVPN SD-WAN Whitepaper." 2023.

        % NIST ABAC 标准文档 (SP 800-162)
        \bibitem{NISTABAC}
        Hu, V. C., Ferraiolo, D., Kuhn, R., Schnitzer, A., Sandlin, K., Miller, R., and Scarfone, K.
        "Guide to Attribute Based Access Control (ABAC) Definition and Considerations."
        \textit{NIST Special Publication 800-162}, 2014.
        Available: \url{https://csrc.nist.gov/pubs/sp/800/162/upd2/final}

        % ZeroTier 平台
        \bibitem{ZeroTier}
        ZeroTier, Inc.
        "ZeroTier: Simple SD-WAN and Zero Trust Access Platform."
        \textit{ZeroTier Official Website}, 2025.
        Available: \url{https://www.zerotier.com/platform/}

        % 贝锐蒲公英 (Oray)
        \bibitem{OrayPgy}
        上海贝锐信息科技股份有限公司.
        "贝锐蒲公英 SD-WAN 解决方案：企业协同办公与远程访问."
        \textit{蒲公英官网}, 2025.
        Available: \url{https://pgy.oray.com/case/scenes/office}

        % EasyTier 开源项目
        \bibitem{EasyTier}
        EasyTier Project.
        "EasyTier: A simple, decentralized mesh VPN with WireGuard support."
        \textit{GitHub Repository}, 2025.
        Available: \url{https://github.com/EasyTier/EasyTier}

    \end{thebibliography}
\end{frame}
\begin{frame}
    \centering
    \Huge \textcolor{blue}{Q \& A} \\
    \vspace{20pt}
    \large 感谢各位老师，请指正。
\end{frame}

\end{document}
